% Save as: module1-problem-set.tex
% ============================================================================
% Problem Set 1: Rational Choice, Preferences & Uncertainty
% Microeconomic Theory for Experimentalists & Applied Microeconomists
% Built from templates/problem_set_template.tex. Compile with pdflatex.
% ============================================================================
\documentclass[11pt]{article}
\usepackage{amsmath, amssymb, amsthm}
\usepackage[margin=1in]{geometry}
\usepackage{enumitem}

\newtheorem{question}{Question}

\title{Problem Set 1: Rational Choice, Preferences \& Uncertainty\\
\large Microeconomic Theory for Experimentalists \& Applied Microeconomists}
\author{Due: \textbf{[date --- before Lecture 1 of Module 2]}}
\date{}

\begin{document}
\maketitle

\noindent\emph{Ground rules: collaboration on ideas is encouraged; write-ups
are individual. You may use AI tools for parts flagged with
$\langle$AI$\rangle$ and must attach the transcript. Show all calculations.}

\medskip

\noindent\textbf{Weights:} Part A --- 30 points. Part B --- 45 points.
Part C --- 25 points.

\bigskip

% ---------------------------------------------------------------------------
\section*{Part A: Theory --- Guided Calculations (30 points)}

\begin{question}[Expected utility with concrete numbers --- 15 points]
An agent has utility $u(x) = \sqrt{x}$ defined over \emph{final wealth}
$x$ (not over gains and losses); initial wealth is zero unless stated
otherwise.

\begin{enumerate}[label=(\alph*)]
  \item \textbf{(4 pts)} Lottery $L$: 50/50 win \$100 or \$0. Compute the
        expected value, expected utility, certainty equivalent, and risk
        premium of $L$. Say in one sentence what the risk premium
        \emph{means} behaviorally.
  \item \textbf{(4 pts)} The agent now has wealth \$100 and faces a loss of
        \$64 with probability $1/4$ (final wealth \$36 if the loss hits).
        Compute her expected utility and certainty equivalent, and show
        that her maximum willingness to pay for full insurance is \$19.
        Compare it to the expected loss and interpret the difference.
  \item \textbf{(4 pts)} An insurer offers full coverage at premium \$18.
        Does she buy? At premium \$20? State the general rule (one
        inequality in words) for when a risk-averse agent buys full
        insurance with a loaded premium.
  \item \textbf{(3 pts)} A classmate concludes from (a) that ``expected
        utility theory predicts people are risk averse.'' Correct them in
        two sentences: what does EU assume, and where does risk aversion
        actually come from in the theory?
\end{enumerate}
\end{question}

\begin{question}[The Allais pattern, by hand --- 15 points]
Consider the two questions from Lecture 1 (Kahneman \& Tversky 1979
payoffs):
\begin{itemize}[nosep]
  \item \textbf{Q1:} $A$ = \$2400 for sure; \quad
        $B$ = 33\% \$2500, 66\% \$2400, 1\% \$0.
  \item \textbf{Q2:} $C$ = 34\% \$2400, 66\% \$0; \quad
        $D$ = 33\% \$2500, 67\% \$0.
\end{itemize}
In KT's data, 82\% chose $A$ and 83\% chose $D$ (the same subjects
answered both questions).

\begin{enumerate}[label=(\alph*)]
  \item \textbf{(6 pts)} Normalize $u(0) = 0$. Write the expected-utility
        inequality for $A \succ B$, simplify it, and show it is the
        \emph{same} inequality as $C \succ D$. Conclude that no utility
        function rationalizes choosing both $A$ and $D$.
        \emph{Hint: cancel the common $0.66\,u(2400)$ term.}
  \item \textbf{(3 pts)} Show that at least 65\% of KT's subjects must have
        chosen the pair $(A, D)$. \emph{Hint: what is the smallest possible
        overlap between the 82\% and the 83\%?}
  \item \textbf{(3 pts)} State in one or two sentences which axiom this
        pattern violates and what the two questions have in common that the
        axiom says must not matter.
  \item \textbf{(3 pts)} A referee writes: ``the Allais pattern is just
        decision noise.'' Give the one-sentence rebuttal that uses the
        \emph{direction} of the violation.
\end{enumerate}
\end{question}

% ---------------------------------------------------------------------------
\section*{Part B: Experimental Design (45 points)}

\begin{question}[Does the certainty effect survive real incentives?]
KT's 82/83 result used hypothetical high stakes. Design an experiment that
tests whether the certainty effect survives \emph{real} incentives, and how
it moves with stake size. Specify:
\begin{enumerate}[label=(\alph*)]
  \item \textbf{(6 pts)} Subject pool and recruitment, and why this
        population is appropriate for the question.
  \item \textbf{(12 pts)} Treatments. Your design must cross incentive
        realness (hypothetical vs.\ paid) with at least two stake levels,
        and you must say what each comparison identifies. Explain how you
        keep the \emph{probabilities and payoff ratios} of the KT questions
        intact while scaling the stakes.
  \item \textbf{(9 pts)} The payment mechanism. If you pay one randomly
        selected decision \emph{per subject}, explain the
        incentive-compatibility argument
        --- and the standard objection to it in \emph{this particular}
        experiment (Lecture 1 flagged it), plus one design choice that
        mitigates the objection.
  \item \textbf{(9 pts)} Hypotheses and rough sample-size reasoning: state
        the violation rate you expect under each treatment (anchor on KT's
        65\% floor and on any published real-stakes replication you know;
        a reasoned guess with a stated source is fine), and the approximate
        $n$ per cell needed to detect the hypothetical-vs-real difference
        you posit. Simulation-based reasoning is welcome; formulas are not
        required.
  \item \textbf{(9 pts)} The biggest confound and your design answer to it.
\end{enumerate}
\end{question}

% ---------------------------------------------------------------------------
\section*{Part C: Silicon Pilot Interpretation $\langle$AI$\rangle$
(25 points)}

\begin{question}[The suspiciously perfect replication]
You run \texttt{module1-simulation-lab.ipynb} with the baseline persona and
the original KT payoffs. Across the 30 sampled responses per cell (at the
lab's fixed temperature), your agents choose $A$ 82\% of the time and $D$
83\% of the time --- matching Kahneman \& Tversky's published shares to the
percentage point. Your lab partner is thrilled: ``the model has human risk
preferences!''

\begin{enumerate}[label=(\alph*)]
  \item \textbf{(8 pts)} Give two mechanisms that could produce this match,
        at least one of which implies your partner is wrong, and say what
        each mechanism predicts the agents would do on a \emph{new} choice
        problem with the same structure but unfamiliar payoffs.
  \item \textbf{(9 pts)} Design the diagnostic simulation that separates
        your two mechanisms using only changes to the prompt or the persona
        grid (e.g.\ novel payoff numbers, disguised framing, paraphrase,
        the ``you are an economist'' persona). State the predicted result
        pattern under each mechanism. Run it if API budget allows and
        report; otherwise state the predictions precisely.
  \item \textbf{(8 pts)} Suppose the diagnostic shows the match is
        recitation, not behavior. Write the two-sentence ``limits of
        silicon subjects'' paragraph you would put in a paper that
        nevertheless uses this lab as a pilot --- what CAN the pilot still
        establish, and what must wait for human subjects?
\end{enumerate}
Attach your AI transcripts and, if you ran code, the results CSV.
\end{question}

\end{document}
